\chapter{Egorov 定理与子列原理}

\begin{solution}{almost-uniform-ae-proof}
取 \(E_k\) 使 \(\mu(E_k)<2^{-k}\) 且补集上一致收敛。由
\(\sum_k\mu(E_k)<\infty\)，集合
\(\limsup_kE_k\) 零测。在其补集上，每个点最终避开全部 \(E_k\)，因此
至少落在某个可任意取大的 \(k\) 的补集内，而该补集上整列收敛。单独一个
\(E_k\) 只有小测度，不一定为零，不能证明几乎处处结论。
\end{solution}

\begin{solution}{egorov-quantifiers}
几乎处处收敛是
\[
 \begin{gathered}
 \exists Z\in\mathcal A,\quad\mu(Z)=0,\quad
 \forall x\in X\setminus Z,\quad\forall\varepsilon>0,\\
 \exists K\ge1,\quad\forall n\ge K:\quad
 |f_n(x)-f(x)|<\varepsilon.
 \end{gathered}
\]
这里的 \(K\) 可依赖 \(x,\varepsilon\)。普通一致收敛不删去点，并要求
对每个 \(\varepsilon>0\)，在选择 \(x\) 以前先取一个 \(K\)，使不等式
对全部 \(x\in X,n\ge K\) 成立。几乎一致收敛是
\[
 \begin{gathered}
 \forall\eta>0,\quad\exists E\in\mathcal A,\quad\mu(E)<\eta,\\
 \forall\varepsilon>0,\quad\exists K\ge1,\quad
 \forall x\in X\setminus E,\quad\forall n\ge K:\\
 |f_n(x)-f(x)|<\varepsilon.
 \end{gathered}
\]
例外集 \(E\) 在选取精度 \(\varepsilon\) 以前已经确定，\(K\) 可依赖
\(\eta,\varepsilon\)，但不依赖余集中的点。Egorov 定理在有限测度下
由第一组量词推出后一组量词；测度可任意小不等于可取零测集。
\end{solution}

\begin{solution}{egorov-xn}
令 \(\delta=\min\{\eta/2,1/2\}\)，取 \(E=(1-\delta,1]\)。则
\(E\subseteq[0,1]\)，\(m(E)=\delta<\eta\)。在 \([0,1-\delta]\) 上，
\[
 \sup x^n=(1-\delta)^n\to0,
\]
故一致收敛。点 \(1\) 本身及其附近是收敛速度不能统一的区域。
\end{solution}

\begin{solution}{egorov-infinite-counterexample}
对每个固定 \(x\)，最终 \(n>x\)，故
\(\mathbf1_{[n,\infty)}(x)=0\)。若在 \(\mathbb R\setminus E\) 上一致
收敛，则存在 \(N\) 使所有 \(n\ge N\) 时函数绝对值小于 \(1/2\)。
特别地 \(n=N\) 给出
\([N,\infty)\cap(\mathbb R\setminus E)=\varnothing\)，所以
\([N,\infty)\subset E\)，其测度无穷，不可能小于给定的有限 \(\eta\)。
\end{solution}

\begin{solution}{borel-cantelli-first}
令 \(F_m=\cup_{n\ge m}E_n\)。则
\[
 \mu(F_m)\le\sum_{n\ge m}\mu(E_n)\longrightarrow0.
\]
而 \(\limsup E_n=\cap_mF_m\subset F_m\) 对所有 \(m\) 成立，所以
\[
 \mu(\limsup E_n)\le\inf_m\mu(F_m)=0.
\]
\end{solution}

\begin{solution}{measure-subsequence-rate}
按依测度收敛递归选 \(n_k\)，使
\[
 \mu(|f_{n_k}-f|>2^{-k})<2^{-k}.
\]
坏集测度之和有限，故几乎每个点只属于有限多个坏集。在其余点上，所给
不等式对所有充分大的 \(k\) 成立。
\end{solution}

\begin{solution}{lp-direct-subsequence}
选 \(n_k\) 使
\(\|f_{n_k}-f\|_p^p<2^{-k}\)。由非负函数逐项积分，
\[
 \int\sum_k|f_{n_k}-f|^p
 =\sum_k\|f_{n_k}-f\|_p^p<\infty.
\]
所以该级数几乎处处有限。收敛级数的项趋零，故
\(|f_{n_k}-f|^p\to0\) 几乎处处。
\end{solution}

\begin{solution}{subsequence-characterization-finite}
正向：每个子列仍依测度收敛，应用子列原理。反向：若不依测度收敛，存在
\(\varepsilon,\eta>0\) 和子列满足
\[
 \mu(|f_{n_k}-f|>\varepsilon)\ge\eta.
\]
假设它有再子列几乎处处收敛于 \(f\)。因 \(\mu(X)<\infty\)，第
\ref{chap:modes-of-convergence} 章的有限测度结论说明该再子列依测度
收敛，故上述超水平集测度应趋零，与下界矛盾。
\end{solution}

\begin{solution}{ae-cauchy-almost-uniform}
先把各 \(f_n\) 取为有限可测代表。其点态 Cauchy 集合为
\[
 C=\bigcap_{r\ge1}\bigcup_{N\ge1}\bigcap_{m,n\ge N}
   \{|f_m-f_n|\le1/r\},
\]
故 \(C\) 可测。若几乎处处 Cauchy，则 \(\mu(X\setminus C)=0\)。
由实数完备性，在 \(C\) 上定义 \(f=\lim_n f_n\)，在 \(X\setminus C\)
上定义 \(f=0\)。它是可测函数列 \(f_n\mathbf1_C\) 的处处有限极限，
所以可测，且 \(f_n\to f\) 几乎处处。
Egorov 定理使它在删去任意小测度集合后一致收敛，因而一致 Cauchy。
反向，取测度小于 \(2^{-k}\) 的例外集 \(E_k\)，使补集上一致 Cauchy。
\(\limsup E_k\) 零测；其补集上的每一点最终落在某个补集 \(E_k^c\) 中，
故数列在该点 Cauchy。
\end{solution}

\begin{solution}{riesz-metric-completeness}
设 \((f_n)\) 在 \(d_\mu\) 下 Cauchy，先选有限可测代表。
由上一章依测度 Cauchy 判据，对每个 \(k\ge1\) 可取整数 \(N_k\)，使
\(m,n\ge N_k\) 时有
\(\mu(\{|f_m-f_n|>2^{-k}\})<2^{-k}\)。把 \(N_k\) 改为递增列，
再依次选严格递增的 \(n_k\ge N_k\)，便有
\[
 \mu(|f_{n_{k+1}}-f_{n_k}|>2^{-k})<2^{-k}.
\]
由第一 Borel--Cantelli 结论，几乎处处有
\(|f_{n_{k+1}}-f_{n_k}|\le2^{-k}\) 最终成立。若一个点从 \(k\ge K\)
起不再落入坏集，则对 \(j>k\ge K\)，
\[
 |f_{n_j}(x)-f_{n_k}(x)|\le\sum_{i=k}^{j-1}2^{-i}
 \le2^{1-k}.
\]
因此子列几乎处处 Cauchy。按上一题构造可测 Cauchy 集合，在其上取有限
极限，在其补集置零，得到有限可测函数 \(f\)。有限测度下 a.e. 收敛蕴含
依测度收敛，故
\(d_\mu(f_{n_k},f)\to0\)。原列 Cauchy，再由三角不等式得
\(d_\mu(f_n,f)\to0\)。具体地，给定 \(\varepsilon>0\)，先取Cauchy尾部
指标，再选尾部内的 \(n_k\) 使 \(d_\mu(f_{n_k},f)<\varepsilon/2\)，
于是尾部每个 \(n\) 都有
\(d_\mu(f_n,f)\le d_\mu(f_n,f_{n_k})+d_\mu(f_{n_k},f)<\varepsilon\)。
\end{solution}

\begin{solution}{egorov-lusin-comparison}
Egorov 从“一个函数列几乎处处收敛”出发，删去小集合后改善为一致收敛；
例外集依赖整个函数列和容许测度 \(\eta\)。Lusin 从“一个可测函数”出发，
本书的版本在有限测度的 Lebesgue 可测集 \(E\subseteq\mathbb R^d\)
中取紧集 \(K\)，使 \(m(E\setminus K)<\eta\)，且 \(f|_K\) 连续。
例外集依赖该函数和 \(\eta\)。

Egorov 证明对每个误差阈值取尾部坏集，利用有限测度得到这些递减坏集
测度趋零，再用可求和的预算合并；它不需要空间带拓扑或测度正则性。
Lusin 证明先用有限测度删去函数值过大的部分，在余集上以有限值简单函数
一致逼近；再由 Lebesgue 内正则性用紧集逼近各水平集。每层有限个不交
紧集的分离性保证简单函数限制连续，最后用连续函数的一致极限传递连续性。
因此两者都删去小集合，但 Lusin 的结论和证明还依赖拓扑，不能只归结为
测度的递减连续性。
\end{solution}
